% ==============================================================
\section{Introduction}\label{sec:intro}
% ==============================================================

Frequent pattern mining has been a cornerstone of data mining research
since the introduction of the Apriori algorithm~\cite{agrawal1994apriori}.
While extensive work has focused on discovering patterns that occur
\emph{frequently} in transactional databases, comparatively little
attention has been paid to the \emph{temporal structure} of pattern
occurrences---specifically, \emph{when} patterns are frequent and
\emph{when} they are not.

Dense interval mining~\cite{dong2012contrast} addresses this gap by
identifying contiguous time periods during which a pattern's sliding-window
support exceeds a given threshold.
These intervals capture epochs of sustained high activity and have
proven useful for event attribution and temporal analysis of purchasing
behavior.
However, the existing framework treats only one side of the coin:
it characterizes \emph{presence} (high support) but provides no
symmetric mechanism for characterizing \emph{absence} (low support).

This asymmetry limits the analytical power of dense interval mining
in several important ways.
First, practitioners cannot directly identify periods when a
previously popular pattern \emph{disappears}---a phenomenon of
significant business interest (e.g., the end of a promotional campaign).
Second, there is no principled framework for comparing how the
temporal structure of a pattern changes across different regimes
(e.g., before vs.\ after an intervention).
Third, the absence of a formal ``anti-dense'' concept prevents the
application of anti-monotonicity-like pruning strategies that could
accelerate the mining of infrequent temporal patterns.

In this paper, we propose two symmetric extensions of the dense
interval mining framework:

\begin{enumerate}
\item \textbf{Anti-Dense Intervals}: We define the dual concept of
dense intervals---maximal contiguous periods where a pattern's support
remains \emph{below} a low threshold $\theta_{\text{low}}$.
Anti-dense intervals capture sustained absence and enable the direct
identification of pattern disappearance epochs.

\item \textbf{Contrast Dense Patterns}: We introduce a framework for
comparing the dense interval \emph{structure} of a pattern across two
temporal regimes.
We define a taxonomy of four canonical topology changes
(Emergence, Vanishing, Amplification, Contraction) and provide a
permutation-based statistical test for assessing whether observed
structural changes are significant.
\end{enumerate}

Our contributions are as follows:
\begin{itemize}
\item We formalize the Anti-Dense Interval and prove a reverse
monotonicity property that enables efficient top-down search
(Section~\ref{sec:problem}).
\item We define the Contrast Dense Pattern framework with a complete
taxonomy of pattern topology changes and a permutation test for
statistical significance (Section~\ref{sec:problem}).
\item We present efficient algorithms that compute anti-dense intervals
and contrast patterns with the same time complexity as existing dense
interval detection (Section~\ref{sec:method}).
\item We validate our approach on synthetic data with known ground truth
and on simulated retail scenarios, demonstrating effective recovery of
anti-dense intervals (F1 = 0.789), accurate classification of topology
changes (76\% accuracy), and statistically significant detection of
campaign-induced pattern vanishing (Section~\ref{sec:experiments}).
\end{itemize}

The remainder of this paper is organized as follows.
Section~\ref{sec:related} reviews related work.
Section~\ref{sec:problem} presents the formal problem definitions.
Section~\ref{sec:method} describes our proposed algorithms.
Section~\ref{sec:experiments} reports experimental results.
Section~\ref{sec:conclusion} concludes with future directions.
